% Подключаем нужные пакеты для таблиц и оформления
\usepackage{booktabs}   % Для красивых горизонтальных линий (\toprule, \midrule, \bottomrule)
\usepackage{multirow}   % Для объединения ячеек по вертикали
\usepackage{caption}    % Для \ContinuedFloat и настройки заголовков
\usepackage{array}      % Для локальной настройки p-колонок таблиц
\usepackage{ragged2e}   % Для аккуратного выравнивания текста в узких колонках
\usepackage{longtable}
\usepackage{etoolbox}

\DeclareCaptionLabelSeparator{emdash}{\space\textemdash\space}  % пробел, длинное тире, пробел

\newcolumntype{P}[1]{>{\RaggedRight\arraybackslash\hspace{0pt}}p{#1}}
\newcommand{\TableHeadTwo}[2]{\shortstack[l]{\textbf{#1}\\\textbf{#2}}}
\newcommand{\TableHeadThree}[3]{\shortstack[l]{\textbf{#1}\\\textbf{#2}\\\textbf{#3}}}

% Настройка заголовков таблиц: слева без абзацного отступа, разделитель — тире (п. 4.6.3)
\captionsetup[table]{position=t, singlelinecheck=false, justification=raggedright, labelsep=emdash, name=Таблица}

% Переключатель для смены выравнивания заголовков при переносе таблицы (п. 4.6.3)
\newcommand*\RightCaption{\captionsetup{justification=raggedleft}}
\newcommand*\LeftCaption{\captionsetup{justification=raggedright}}

% ————————
% Шаблон: Многоуровневый заголовок
% ————————
\newcommand{\ExampleMultiLevelTable}{%
\begin{table}[htbp]
\caption{Пример построения таблицы}
\centering
\begin{tabular}{|c|c|c|c|}
\hline
\multirow{2}{*}{\textbf{Заголовок}} & \multicolumn{2}{c|}{\textbf{Заголовок колонки}} & \multirow{2}{*}{\textbf{Заголовок колонки}}\\
\cline{2-3}
& \textbf{Подзаголовок} & \textbf{Подзаголовок} & \\
\hline
A & B & C & D \\
\hline
\end{tabular}
\end{table}
}

% Шаблон: «Разбитая» таблица (две части под одним номером)
% Номер таблицы в строке «Продолжение таблицы N» подставляется автоматически
\newcommand{\ExampleSplitTable}{%
  % Первая часть
  \begin{table}[htbp]
    \LeftCaption
    \caption{Пример переноса таблицы}
    \label{tab:split_example}
    \centering
    \begin{tabular}{|c|c|c|c|}
    \hline
    \textbf{Заголовок} & \textbf{Заголовок колонки} & \textbf{Заголовок колонки} & \textbf{Заголовок колонки} \\
    \hline
    1 & 2 & 3 & 4 \\
    \hline
    2 & 3 & 4 & 5 \\
    \hline
    \end{tabular}
  \end{table}

  % Продолжение: номер подставляется автоматически через \thetable
  \begin{table}[htbp]
    \ContinuedFloat
    \RightCaption
    \caption*{Продолжение таблицы \thetable}
    \centering
    \begin{tabular}{|c|c|c|c|}
    \hline
    \textbf{Заголовок} & \textbf{Заголовок колонки} & \textbf{Заголовок колонки} & \textbf{Заголовок колонки} \\
    \hline
    3 & 4 & 5 & 6 \\
    \hline
    4 & 5 & 6 & 7 \\
    \hline
    \end{tabular}
  \end{table}
}
